% Options for packages loaded elsewhere
\PassOptionsToPackage{unicode}{hyperref}
\PassOptionsToPackage{hyphens}{url}
\documentclass[
]{article}
\usepackage{xcolor}
\usepackage[margin=1in]{geometry}
\usepackage{amsmath,amssymb}
\setcounter{secnumdepth}{5}
\usepackage{iftex}
\ifPDFTeX
  \usepackage[T1]{fontenc}
  \usepackage[utf8]{inputenc}
  \usepackage{textcomp} % provide euro and other symbols
\else % if luatex or xetex
  \usepackage{unicode-math} % this also loads fontspec
  \defaultfontfeatures{Scale=MatchLowercase}
  \defaultfontfeatures[\rmfamily]{Ligatures=TeX,Scale=1}
\fi
\usepackage{lmodern}
\ifPDFTeX\else
  % xetex/luatex font selection
\fi
% Use upquote if available, for straight quotes in verbatim environments
\IfFileExists{upquote.sty}{\usepackage{upquote}}{}
\IfFileExists{microtype.sty}{% use microtype if available
  \usepackage[]{microtype}
  \UseMicrotypeSet[protrusion]{basicmath} % disable protrusion for tt fonts
}{}
\makeatletter
\@ifundefined{KOMAClassName}{% if non-KOMA class
  \IfFileExists{parskip.sty}{%
    \usepackage{parskip}
  }{% else
    \setlength{\parindent}{0pt}
    \setlength{\parskip}{6pt plus 2pt minus 1pt}}
}{% if KOMA class
  \KOMAoptions{parskip=half}}
\makeatother
\usepackage{graphicx}
\makeatletter
\newsavebox\pandoc@box
\newcommand*\pandocbounded[1]{% scales image to fit in text height/width
  \sbox\pandoc@box{#1}%
  \Gscale@div\@tempa{\textheight}{\dimexpr\ht\pandoc@box+\dp\pandoc@box\relax}%
  \Gscale@div\@tempb{\linewidth}{\wd\pandoc@box}%
  \ifdim\@tempb\p@<\@tempa\p@\let\@tempa\@tempb\fi% select the smaller of both
  \ifdim\@tempa\p@<\p@\scalebox{\@tempa}{\usebox\pandoc@box}%
  \else\usebox{\pandoc@box}%
  \fi%
}
% Set default figure placement to htbp
\def\fps@figure{htbp}
\makeatother
\setlength{\emergencystretch}{3em} % prevent overfull lines
\providecommand{\tightlist}{%
  \setlength{\itemsep}{0pt}\setlength{\parskip}{0pt}}
\usepackage{amsmath}
\usepackage{amssymb}
\usepackage{amsthm}
\usepackage{geometry}
\geometry{margin=1in}
\usepackage{bookmark}
\IfFileExists{xurl.sty}{\usepackage{xurl}}{} % add URL line breaks if available
\urlstyle{same}
\hypersetup{
  pdftitle={Derivações Detalhadas das Proposições 1 e 2 (IVB)},
  pdfauthor={Manoel Galdino},
  hidelinks,
  pdfcreator={LaTeX via pandoc}}

\title{Derivações Detalhadas das Proposições 1 e 2 (IVB)}
\author{Manoel Galdino}
\date{27/02/2026}

\begin{document}
\maketitle

\section{Q\&A Inicial}\label{qa-inicial}

\textbf{Q1. Qual é o objetivo deste documento?}\\
Checar, passo a passo, as derivações da Proposição 1 (cross-section) e
da Proposição 2 (ADL(1,0)) do IVB, com expansão explícita dos pontos que
geraram dúvidas.

\textbf{Q2. Quais foram os pontos de atenção incorporados aqui?}\\
(i) passagem da condição normal para a fórmula de coeficiente,\\
(ii) expansão da covariância com médias explícitas,\\
(iii) cuidado com sinais algébricos do tipo \(- (a-b)\),\\
(iv) residualização via FWL,\\
(v) diferença entre ortogonalidade OLS e ausência de viés causal.

\textbf{Q3. Qual notação será usada para facilitar a leitura?}\\
Usaremos: \[
\text{(M1)} = \text{modelo curto (correto)}, \qquad
\text{(M2)} = \text{modelo longo (inclui collider)}.
\]

\section{Premissas Mínimas}\label{premissas-muxednimas}

Para as razões com variância no denominador estarem bem definidas: \[
\operatorname{Var}(D)>0
\] e, no caso dinâmico, \[
\operatorname{Var}(\widetilde{D}_t)>0.
\] Assumimos também momentos de segunda ordem finitos.

\section{Proposição 1: Caso
Cross-Section}\label{proposiuxe7uxe3o-1-caso-cross-section}

\subsection{Enunciado}\label{enunciado}

Quando um collider \(Z\) é incluído por engano na regressão de \(y\) em
\(D\), \[
\beta_1^\star - \beta_1 = -\beta_2^\star \phi_1,
\] onde \(\beta_2^\star\) é o coeficiente de \(Z\) na regressão longa e
\(\phi_1\) é o coeficiente da regressão auxiliar de \(Z\) em \(D\).

\subsection{Passo 1. Escrever os dois
modelos}\label{passo-1.-escrever-os-dois-modelos}

\[
\text{(M1)} \quad y = \beta_0 + \beta_1 D + e,
\] \[
\text{(M2)} \quad y = \beta_0^\star + \beta_1^\star D + \beta_2^\star Z + e^\star.
\]

No contexto do paper, (M1) é o modelo causal correto e (M2) é o modelo
misspecificado por incluir o collider.

\subsection{\texorpdfstring{Passo 2. Coeficiente de \(D\) no modelo
curto}{Passo 2. Coeficiente de D no modelo curto}}\label{passo-2.-coeficiente-de-d-no-modelo-curto}

\[
\beta_1 = \frac{\operatorname{Cov}(D,y)}{\operatorname{Var}(D)}.
\]

\subsection{Passo 3. Condição normal no modelo longo (FOC
explícita)}\label{passo-3.-condiuxe7uxe3o-normal-no-modelo-longo-foc-expluxedcita}

Defina o resíduo no modelo longo: \[
u^\star = y - \beta_0^\star - \beta_1^\star D - \beta_2^\star Z.
\]

OLS populacional escolhe \((\beta_0^\star,\beta_1^\star,\beta_2^\star)\)
para minimizar: \[
Q = \mathbb{E}\!\left[(u^\star)^2\right].
\]

Derivando em relação a \(\beta_1^\star\): \[
\frac{\partial Q}{\partial \beta_1^\star}
=
\mathbb{E}\!\left[2u^\star\left(-D\right)\right]
=
-2\mathbb{E}[Du^\star]=0
\;\Rightarrow\;
\mathbb{E}[Du^\star]=0.
\]

Com intercepto, também vale \(\mathbb{E}[u^\star]=0\). Logo: \[
\operatorname{Cov}(D,u^\star)=\mathbb{E}[Du^\star]-\mathbb{E}[D]\mathbb{E}[u^\star]=0.
\] Isto implica: \[
\operatorname{Cov}\!\big(D,\;y-\beta_0^\star-\beta_1^\star D-\beta_2^\star Z\big)=0.
\]

\subsection{Passo 4. Expansão completa da covariância (com médias
explícitas)}\label{passo-4.-expansuxe3o-completa-da-covariuxe2ncia-com-muxe9dias-expluxedcitas}

Partindo de: \[
0=\operatorname{Cov}\!\big(D,\;y-\beta_0^\star-\beta_1^\star D-\beta_2^\star Z\big),
\] aplique a definição: \[
\operatorname{Cov}(A,B)=\mathbb{E}\!\left[(A-\mu_A)(B-\mu_B)\right].
\]

Assim: \[
0=\mathbb{E}\!\left[(D-\mu_D)\left(\Xi-\mu_\Xi\right)\right],
\] onde \[
\Xi=y-\beta_0^\star-\beta_1^\star D-\beta_2^\star Z,
\quad
\mu_\Xi=\mu_y-\beta_0^\star-\beta_1^\star\mu_D-\beta_2^\star\mu_Z.
\]

Logo: \[
\Xi-\mu_\Xi
=(y-\mu_y)-\beta_1^\star(D-\mu_D)-\beta_2^\star(Z-\mu_Z).
\]

Substituindo: \begin{align*}
0
&=
\mathbb{E}\!\left[(D-\mu_D)(y-\mu_y)\right]
-\beta_1^\star\mathbb{E}\!\left[(D-\mu_D)^2\right]
-\beta_2^\star\mathbb{E}\!\left[(D-\mu_D)(Z-\mu_Z)\right] \\
&=
\operatorname{Cov}(D,y)-\beta_1^\star\operatorname{Var}(D)-\beta_2^\star\operatorname{Cov}(D,Z).
\end{align*}

Isolando \(\beta_1^\star\): \[
\beta_1^\star
=
\frac{\operatorname{Cov}(D,y)-\beta_2^\star\operatorname{Cov}(D,Z)}
{\operatorname{Var}(D)}.
\]

\subsection{\texorpdfstring{Passo 5. Subtrair
\(\beta_1\)}{Passo 5. Subtrair \textbackslash beta\_1}}\label{passo-5.-subtrair-beta_1}

\[
\beta_1^\star-\beta_1
=
-\beta_2^\star\frac{\operatorname{Cov}(D,Z)}{\operatorname{Var}(D)}.
\]

\subsection{\texorpdfstring{Passo 6. Regressão auxiliar para
\(Z\)}{Passo 6. Regressão auxiliar para Z}}\label{passo-6.-regressuxe3o-auxiliar-para-z}

\[
Z=\phi_0+\phi_1 D+\varepsilon
\quad\Rightarrow\quad
\phi_1=\frac{\operatorname{Cov}(D,Z)}{\operatorname{Var}(D)}.
\]

Substituindo: \[
\boxed{\beta_1^\star-\beta_1=-\beta_2^\star\phi_1.}
\]

\subsection{Nota Didática A: Sinal
Algébrico}\label{nota-diduxe1tica-a-sinal-alguxe9brico}

Um ponto crítico de sinal: \[
-(a-b)=-a+b.
\] Portanto, \[
-\beta_1^\star D -(-\beta_1^\star\mu_D)
=
-\beta_1^\star D+\beta_1^\star\mu_D
=
-\beta_1^\star(D-\mu_D),
\] e não \(-\beta_1^\star(D+\mu_D)\).

\section{Proposição 2: Caso ADL(1,0) via
FWL}\label{proposiuxe7uxe3o-2-caso-adl10-via-fwl}

\subsection{Enunciado}\label{enunciado-1}

No ADL(1,0), ao incluir um collider \(Z_t\): \[
\beta^\star-\beta=-\theta^\star\pi.
\]

\subsection{Passo 1. Modelos (curto e
longo)}\label{passo-1.-modelos-curto-e-longo}

\[
\text{(M1)}\quad y_t=\alpha+\beta D_t+\rho y_{t-1}+e_t,
\] \[
\text{(M2)}\quad y_t=\alpha^\star+\beta^\star D_t+\rho^\star y_{t-1}+\theta^\star Z_t+v_t.
\]

Aqui, o objeto de viés \(\beta^\star-\beta\) compara duas regressões que
diferem apenas pela inclusão de \(Z_t\).

\subsection{Passo 2. Residualização explícita
(FWL)}\label{passo-2.-residualizauxe7uxe3o-expluxedcita-fwl}

Defina controles legítimos: \[
W_t=(1,y_{t-1}).
\]

Residualizar significa rodar: \[
y_t = a_y + b_y y_{t-1} + r^y_t,\quad \widetilde{y}_t=r^y_t,
\] \[
D_t = a_D + b_D y_{t-1} + r^D_t,\quad \widetilde{D}_t=r^D_t,
\] \[
Z_t = a_Z + b_Z y_{t-1} + r^Z_t,\quad \widetilde{Z}_t=r^Z_t.
\]

Pelo FWL, a regressão longa equivale a: \[
\widetilde{y}_t=\beta^\star\widetilde{D}_t+\theta^\star\widetilde{Z}_t+\widetilde{v}_t.
\]

\subsection{\texorpdfstring{Passo 3. Da minimização à FOC para
\(\beta^\star\)}{Passo 3. Da minimização à FOC para \textbackslash beta\^{}\textbackslash star}}\label{passo-3.-da-minimizauxe7uxe3o-uxe0-foc-para-betastar}

Considere: \[
Q(\beta^\star,\theta^\star)=
\mathbb{E}\!\left[\left(\widetilde{y}_t-\beta^\star\widetilde{D}_t-\theta^\star\widetilde{Z}_t\right)^2\right].
\]

Derivando em relação a \(\beta^\star\): \[
\frac{\partial Q}{\partial \beta^\star}
=
-2\,\mathbb{E}\!\left[\widetilde{D}_t\left(\widetilde{y}_t-\beta^\star\widetilde{D}_t-\theta^\star\widetilde{Z}_t\right)\right]=0.
\]

Então: \[
\mathbb{E}[\widetilde{D}_t\widetilde{y}_t]
-\beta^\star\mathbb{E}[\widetilde{D}_t^2]
-\theta^\star\mathbb{E}[\widetilde{D}_t\widetilde{Z}_t]=0.
\]

\subsection{\texorpdfstring{Passo 4. De \(\mathbb{E}[\cdot]\) para
covariância/variância}{Passo 4. De \textbackslash mathbb\{E\}{[}\textbackslash cdot{]} para covariância/variância}}\label{passo-4.-de-mathbbecdot-para-covariuxe2nciavariuxe2ncia}

Como há intercepto em \(W_t\), os resíduos têm média zero: \[
\mathbb{E}[\widetilde{D}_t]=\mathbb{E}[\widetilde{y}_t]=\mathbb{E}[\widetilde{Z}_t]=0.
\]

Logo: \[
\mathbb{E}[\widetilde{D}_t\widetilde{y}_t]
=
\operatorname{Cov}(\widetilde{D}_t,\widetilde{y}_t),
\quad
\mathbb{E}[\widetilde{D}_t^2]
=
\operatorname{Var}(\widetilde{D}_t),
\] \[
\mathbb{E}[\widetilde{D}_t\widetilde{Z}_t]
=
\operatorname{Cov}(\widetilde{D}_t,\widetilde{Z}_t).
\]

Portanto: \[
\operatorname{Cov}(\widetilde{D}_t,\widetilde{y}_t)
-\beta^\star\operatorname{Var}(\widetilde{D}_t)
-\theta^\star\operatorname{Cov}(\widetilde{D}_t,\widetilde{Z}_t)=0.
\]

Isolando \(\beta^\star\): \[
\beta^\star=
\frac{\operatorname{Cov}(\widetilde{D}_t,\widetilde{y}_t)-\theta^\star\operatorname{Cov}(\widetilde{D}_t,\widetilde{Z}_t)}
{\operatorname{Var}(\widetilde{D}_t)}.
\]

\subsection{\texorpdfstring{Passo 5. Repetir no modelo curto para obter
\(\beta\)}{Passo 5. Repetir no modelo curto para obter \textbackslash beta}}\label{passo-5.-repetir-no-modelo-curto-para-obter-beta}

No curto residualizado: \[
\widetilde{y}_t=\beta\widetilde{D}_t+\widetilde{e}_t.
\]

Pela mesma lógica: \[
\beta=\frac{\operatorname{Cov}(\widetilde{D}_t,\widetilde{y}_t)}
{\operatorname{Var}(\widetilde{D}_t)}.
\]

\subsection{Passo 6. Subtração (diferença entre longo e
curto)}\label{passo-6.-subtrauxe7uxe3o-diferenuxe7a-entre-longo-e-curto}

\[
\beta^\star-\beta
=
-\theta^\star
\frac{\operatorname{Cov}(\widetilde{D}_t,\widetilde{Z}_t)}
{\operatorname{Var}(\widetilde{D}_t)}.
\]

\subsection{\texorpdfstring{Passo 7. Regressão auxiliar e definição de
\(\pi\)}{Passo 7. Regressão auxiliar e definição de \textbackslash pi}}\label{passo-7.-regressuxe3o-auxiliar-e-definiuxe7uxe3o-de-pi}

\[
Z_t=a+\pi D_t+\lambda y_{t-1}+\eta_t.
\] No espaço residualizado por \(W_t\): \[
\widetilde{Z}_t=\pi\widetilde{D}_t+\widetilde{\eta}_t
\quad\Rightarrow\quad
\pi=
\frac{\operatorname{Cov}(\widetilde{D}_t,\widetilde{Z}_t)}
{\operatorname{Var}(\widetilde{D}_t)}.
\]

Substituindo: \[
\boxed{\beta^\star-\beta=-\theta^\star\pi.}
\]

\subsection{Nota Didática B: Ortogonalidade OLS vs Viés
Causal}\label{nota-diduxe1tica-b-ortogonalidade-ols-vs-viuxe9s-causal}

No longo residualizado, vale por construção OLS: \[
\mathbb{E}[\widetilde{D}_t\widetilde{v}_t]=0.
\] Isso \textbf{não} implica \(\beta^\star=\beta\).\\
Implica apenas que \(\beta^\star\) é o coeficiente de projeção linear do
modelo longo.\\
O viés causal em relação ao alvo do modelo curto continua sendo: \[
\beta^\star-\beta.
\]

\subsection{\texorpdfstring{Nota Didática C: Viés em \(M2\) vem de
\(Z_t\) ou de
\(y_{t-1}\)?}{Nota Didática C: Viés em M2 vem de Z\_t ou de y\_\{t-1\}?}}\label{nota-diduxe1tica-c-viuxe9s-em-m2-vem-de-z_t-ou-de-y_t-1}

Na decomposição \(\beta^\star-\beta\), o termo capturado é o efeito de
incluir \(Z_t\), pois \(M1\) e \(M2\) compartilham \(y_{t-1}\).\\
Ou seja, a diferença entre os modelos é apenas a presença do collider.

\section{Conclusão}\label{conclusuxe3o}

As duas proposições têm a mesma estrutura: \[
\text{IVB}=-\theta^\star\times\pi.
\] No caso cross-section,
\((\theta^\star,\pi)=(\beta_2^\star,\phi_1)\).\\
No caso ADL(1,0), ambos são efeitos parciais no espaço residualizado por
\(W_t\).

O ponto central é que a identidade é algébrica e operacional: os
componentes são estimáveis por regressões padrão, mas a interpretação
causal depende de justificar substantivamente que \(Z\) é collider.

\end{document}
